
@@M0@@
Or @@M1@@ d'après (5.2.3).
Utilisant la relation (1.12.1), on obtient finalement :
@@M2@@
d'où les relations cherchées

\textbf{Proposition 5.10.} \uline{Avec les notations et les hypothèses de} 5.1, \uline{soit} @@M3@@. \uline{Pour que} @@M4@@, \uline{il faut et il suffit qu'il existe} @@M5@@, \uline{tel que} :
@@M6@@
De @@M7@@, on déduit @@M8@@. La condition est donc suffisante.

Soit @@M9@@ tel que @@M10@@, soit @@M11@@. D'après la remarque 1.13, on peut, de façon unique, trouver des @@M12@@ tels que :
@@M13@@
Par ailleurs, les @@M14@@ vérifient une équation de dépendance :
@@M15@@
Comme @@M16@@, on obtient :
@@M17@@
Utilisant l'indépendance de @@M18@@, on trouve :
@@M19@@
